\chapter{Conclusion and Future Work}
\label{ch:conclusion}

\section{Conclusion}

This project set out to allocate tasks in a rover fleet whose communication
graph changes every few seconds, using a constrained QUBO to prune the graph
and a trust-aware GNN to assign tasks on what remains. It delivered that
pipeline, and it delivered something more valuable than the pipeline alone:
a careful account of what each part contributes, measured against the
strongest alternatives available.

For the optimisation part, we derived the QUBO encoding and checked it
against an independent re-derivation and against exhaustive enumeration. A
five-rover example showed that, because of how it is built, the encoding
rejects the connectivity traps that a pruner using only soft scores would
fall into. The solver audit then showed that a correct encoding is not
enough on its own. Four well-known annealing-type samplers found the exact
optimum on only 48 to 64\% of the instances we could verify. The penalty
walls that make a violation unprofitable also make it very hard for these
samplers to cross from one feasible state to another. Once the search was
restricted to feasible states, this barrier disappeared. The projected
annealer is exact on all of the verifiable instances and stays within
$1.4\times10^{-3}$ of the exact optimum on all of the larger ones. Last, the
exact MILP baseline is exact as well and about 100 times faster. For this
reason we see the QUBO as an encoding for quantum hardware and not as the
best way to solve the problem on a classical computer.

For the learning part, trust-aware assignment raises mission success by 8.1
points compared with a trust-only baseline, and the confidence interval for
this gain does not include zero. Against random assignment the gain is 13.8
points. These results come from a simulator in which success depends causally
on both the assignment and the pruned topology. Message passing adds a small margin that is not yet significant,
and a cloud-trained readout did not beat the fixed weights on the data
available.

On the systems side, the pipeline runs unchanged in a containerised Gazebo
Harmonic and ROS~2 Jazzy world against live odometry, and in an interactive
3D engine whose rovers drive with exact, provably collision-free
differential-drive kinematics, whose renderer costs a constant seven draw
calls, and whose comparisons against weaker baselines are paired and
statistically annotated.

\section{Key Take-Aways}

\begin{enumerate}[label=\arabic*.]
\item \textbf{Bit-level annealers may fail even on a correct QUBO.} Before
drawing any conclusion about a formulation, the solver's answers should be
compared with an exact solution.
\item \textbf{On classical hardware, solve the constrained problem directly.}
The MILP gives the exact answer and is two orders of magnitude faster. The
QUBO is only worth using as an input format for annealing hardware.
\item \textbf{The guarantees are not free.} The constrained pruner is the only
arm that never breaks a constraint, but at ten rovers it loses 2 to 4 points
of raw success, and none of the pressure levels in our sweep turned the
guarantee into higher success. It makes sense to use it when breaking a
constraint is not acceptable at all.
\item \textbf{Most of the assignment gain comes from trust.} The trust model,
the gates and delivery awareness give most of the improvement, and message
passing adds only a small amount on top.
\item \textbf{Audit the audit.} Logging which constraints were actually
enforced exposed a silently relaxed budget in the project's own earlier
results, and changed its headline table.
\end{enumerate}

\section{Performance Against the Objectives}

Table~\ref{tab:conc:objectives} compares the outcomes with the objectives set
in Chapter~\ref{ch:intro}.

\begin{table}[htbp]
\centering
\caption{Outcome of each objective.}
\label{tab:conc:objectives}
\small
\begin{tabularx}{\linewidth}{@{}l X@{}}
\toprule
Objective & Outcome \\
\midrule
1. Constrained QUBO & Achieved: utility, stability, degree, budget and exact flow-based connectivity in one energy \\
2. Verified encoding & Achieved: agreement within $2.3\times10^{-10}$; no violated exact optimum in 38 instances \\
3. Reliable solver & Achieved with a caveat: projected annealer exact on all verifiable instances; generic samplers are not \\
4. Exact baselines & Achieved: direct MILP and exact linearised solve; MILP is faster than any annealer \\
5. Trust-aware GNN & Achieved: $+8.1$ pts over trust-only; message-passing margin and trained readout not significant \\
6. Causal evaluation & Achieved: causal delivery model, repaired baselines, pressure sweep, paired statistics \\
7. Deployments & Achieved: Gazebo and ROS~2 integration; interactive 3D engine with live comparisons \\
\bottomrule
\end{tabularx}
\end{table}

It is also worth being clear about what we built ourselves and what we
used from others. We built the QUBO formulation and its verification, the
projected annealer, the causal simulator, the evaluation protocol, the
planner node, the planner service and the 3D engine. The samplers, the CBC
solver, Gazebo, ROS~2 and the rendering libraries were used as published.

\section{Benefits}

Apart from the pipeline itself, the work is useful in three ways. The first
is the way it evaluates things: exact referees, repaired baselines, causal
metrics, paired statistics and a record of which constraints were relaxed can
be reused for any QUBO-based graph method. The second is a hardware experiment
that is ready to run, because the exported Q matrices and their certified
optima form a benchmark on which quantum annealers can be compared fairly. The
third is a tool for teaching and demonstrations, since the 3D
engine lets people see constrained optimisation, trust learning and the
trade-offs between them without having to read through tables of numbers.

\section{Future Work}

\begin{enumerate}[label=\arabic*.]
\item \textbf{Moving rovers in Gazebo.} Publish velocity commands from the
planner node, using the controller of the 3D engine, so that the candidate
graph changes in the physics simulation as it does in the browser.
\item \textbf{Physical validation.} Build the four-rover ESP32 testbed and
replace the synthetic reliability and radio models with measured ones.
\item \textbf{Quantum annealing hardware.} The exported Q matrices could be
run on a real quantum annealer and the results compared with the certified
optima. This would show whether transverse-field dynamics handle the penalty
barrier better than the classical samplers did.
\item \textbf{Delivery-aware utility.} The pruner's utility ranks links by
their own quality and misses links that delivery depends on; adding
path-criticality to the objective addresses the 0.67 recall directly.
\item \textbf{Scalable encodings.} To handle hundreds of rovers, the dense
matrix and the flow variables would need to be replaced by a sparse or
hierarchical encoding.
\item \textbf{Richer trust and on-policy training.} Contextual trust
estimates and a larger exploration budget for training the readout.
\item \textbf{Statistical power.} More seeds and scenarios to resolve the
small margins, such as message passing over the no-graph readout.
\end{enumerate}
